\section{Inverse-color exposure and balanced trace moments}\label{tra:inverse-color}

The density-cap compatibility bound below is a consequence of the
weighted theorem in Section~\ref{tra:duality}. We give a different proof
by exposing inverse colors. It compares the laws of remaining positions
through their original-column labels, charges large atoms directly to
column entropy, and handles repeated labels by a separate exponential
moment. The row array is first fixed; within that probability calculation,
inverse permutations are exposed by destination color rather than by
column prefixes.

The operator application uses spectral projection ranks. Its separate
sparse input is an exact endpoint-kernel product formula, which gives
a forest estimate through first-row level $n^{.64}$.

Put $P_n=T_nT_n^*$ and define $Z_n(s)=\operatorname{Tr}P_n^s$ for $s>0$. Thus
$Z_n(s)=\operatorname{Tr}Q_n^s$; the two positive squares have the
same eigenvalues, although they are generally different operators.
Traces retain regular multiplicities.

\subsection{A compatibility bound from inverse colors}
\begin{proposition}\label{tra:inverse-color:compatibility}
\label{tra:inverse-color:compatibility-statement}
Let \(n=qm\), with \(m\le q\le2m\) and \(m\) sufficiently large. On a \(q\times m\) grid choose all row and column permutations independently. Relative to uniform in each line, suppose their laws have bounded densities, and let \(L_R,L_C\) be the sums for rows and columns, respectively, of the logarithms of these density bounds.

Consider the product of the row permutation layer and then the column permutation layer. Call the product reversible in order if it can also be written as a column layer followed by a row layer (this just refers to a factorization of a permutation). With \(C\) an absolute constant,
\begin{equation}
 \Pr(\text{reversible in order})
 \ \le\
 \exp\left(-n+n^{0.58}+\frac{C}{\log n}(L_R+L_C)\right).                         \label{tra:inverse-color:eq1}
\end{equation}
\end{proposition}
\begin{proof}\label{tra:inverse-color:compatibility-proof}
\label{tra:inverse-color:inverse-column-entropy-and-repeat-moment}
To see what the condition means, write the row permutations as \(a_i\), column permutations as \(b_j\), and at each intermediate cell \((i,j)\) put
\[
                    X_{ij}=a_i^{-1}(j),\qquad Y_{ij}=b_j(i).
\]
Then the condition is precisely that the \(n\) pairs \((X_{ij},Y_{ij})\) are all distinct. In fact in a column-first factorization, all cards from any one original column must have different destination rows, and this is also sufficient, since then one can first put all cards in the correct row by column permutations, before going to their destinations by row permutations.

Fix the array \(X\). For each \(j,x\) let \(M_{jx}=|\{i:X_{ij}=x\}|\), and let
\[
                   W=\sum_{j,x:\ M_{jx}>m^{0.05}} M_{jx}.
\]
With only column randomness now, let \(p_{\rm ok}\) be the probability of the condition. We claim
\begin{equation}
                \log p_{\rm ok}
                  \le -n+W+n^{0.57}+\frac{C}{\log m}L_C .                       \label{tra:inverse-color:eq2}
\end{equation}
Assume the probability is positive, and let \(\nu\) be the conditional law of the column permutations given success. Denote column laws before conditioning by \(\mu_j\), and the marginals of \(\nu\) by \(\nu_j\). With relative entropy denoted by \(D_{\rm KL}\), put
\[
         I=D_{\rm KL}(\nu\ \|\ \textstyle\prod_j\nu_j),\qquad
         D=\sum_j D_{\rm KL}(\nu_j\ \|\ \mu_j).
\]
We have
\begin{equation}
 -\log p_{\rm ok}=I+D,\qquad
 \sum_j D_{\rm KL}(\nu_j\ \|\ \text{uniform})\le D+L_C .                       \label{tra:inverse-color:eq3}
\end{equation}

Estimate \(I\) by revealing inverse permutations \(i_j(y)=b_j^{-1}(y)\) in the order of colors \(y=1,\ldots,q\) (destination row is the color), and within each color in uniformly random order of columns, independent of the configuration. For each color we can order by independent uniform priorities in \([0,1]\), low priority number revealed first. In the chain rule for \(I\), consider the divergence contribution when \(i_j(y)\) is revealed. Under the reference product of the marginals, the prediction law \(u_i\) on possible \(i\)'s for \(i_j(y)\) is determined by the previously revealed colors in that column only. Let
\[
                          v_x=\sum_{i:X_{ij}=x} u_i.
\]
Under the true conditional reveal law from \(\nu\), the \(x\) value this color gets from column \(j\) cannot have already been used by this color in another column. So the KL contribution (conditional divergence) is at least minus the log of the total \(v\)-mass of \(x\)'s not yet used by the color.

We average this last bound. Fix a full sample from \(\nu\), and for column \(j\) in that color fix its priority \(t\); the vector \(v\) does not depend on the other priorities. Every \(x\) except the actual value \(x_* = X_{i_j(y),j}\) of the sample at this reveal is used by the color in another column, hence has probability \(t\) of being used earlier. By Jensen the expected lower bound on divergence from the allowed mass is at least
\[
                              -\log(1-t+t v_{x_*}).
\]
Integrating over \(t\), this is
\[
                 1-h(v_{x_*}),\qquad h(z)=\frac{-z\log z}{1-z},
\]
using continuous limit values. Here \(0\le h(z)\le 1\), and \(h(z)\le C z\log(e/z)\) on \([0,1]\).

For the whole sum of errors \(h(v_{x_*})\), the terms with \(M_{j x_*}>m^{0.05}\) cost at most \(W\) since all cells in each column are revealed. To bound the other errors, for column \(j\) and color \(y\) condition just on the previous colors in that column. Then under sampling from \(\nu\) the conditional law of \(x_*\) is \(v\). Write \(R=q-y+1\) for the number of remaining rows. For \(M_{jx}\le m^{0.05}\):
\begin{itemize}
\item Those \(x\) with \(v_x\le m^{0.1}/R\) have \(h(v_x)\le C(\log m)m^{0.1}/R\), using also \(R\le 2m\).
\item The total \(v\)-probability of the others with \(M_{jx}\le m^{0.05}\) is at most
\end{itemize}
\[
                          \frac{C}{\log m}D_{\rm KL}(u\ \|\ U_R),
\]
where \(U_R\) is uniform among the remaining \(i\)'s in this column. Indeed under \(U_R\), the induced probability of any such \(x\) is at most \(m^{0.05}/R\); compare this induced law with \(v\), writing the divergence as a sum of nonnegative terms \(a\log(a/b)-a+b\), and use the data processing inequality. On the \(x\)'s in question we have probability ratio at least \(m^{0.05}\).

Sum the expected errors. For the first bullet the summed bound is at most \(C m^{1.1}(\log m)^2\le n^{0.57}\) for sufficiently large \(m\). For the second bullet we can use \(h\le1\); the divergences in the bound sum (in expectation) to the left side of the second inequality in \eqref{tra:inverse-color:eq3}, by the chain rule for the column marginals. We conclude
\[
                 I\ge n-W-n^{0.57}-\frac{C}{\log m}(D+L_C).
\]
Combining with \eqref{tra:inverse-color:eq3} proves \eqref{tra:inverse-color:eq2}, since \(m\) is sufficiently large.

The inverse-color exposure has reduced the loss to repeated original
column labels. We now average over the row permutations by bounding
the exponential cost of $W$. First under uniform row permutations,
\begin{equation}
                        {\mathbb E}_{\rm unif}\exp(c(\log m)W)\le 2              \label{tra:inverse-color:eq4}
\end{equation}
for some absolute \(c>0\) and sufficiently large \(m\). Here are details of the count. Groups of entries contributing to \(W\) are specified by a pair \(j,x\), a size \(l>m^{0.05}\), and \(l\) chosen rows where \(X_{ij}=x\). For any collection with distinct indexing pairs \(j,x\), fulfilling its requirements has probability at most \((e/m)^{\sum l}\) (or zero if requirements conflict). In fact if \(b\) entries are prescribed in a row their probability is at most \(1/(m)_b\le(e/m)^b\), with falling factorial notation. Sum over collections with the exponential weight for their total size; this bounds the moment because we can take the collection of all actual contributing groups, and include the empty collection. We obtain the upper bound
\[
                 \left(1+\sum_{l>m^{0.05}}^q
                   \binom ql(e/m)^l m^{cl}\right)^{m^2},
\]
where the sum uses integral \(l\). By \(\binom ql\le(2em/l)^l\), say \(c=0.01\) suffices for \eqref{tra:inverse-color:eq4}. Changing the row laws to the ones in \eqref{tra:inverse-color:eq1}, the joint density of rows is at most \(\exp(L_R)\). Thus by Hölder, putting \(B=c\log m>1\),
\[
                     \mathbb E\exp W\le (2\exp L_R)^{1/B}.
\]
Averaging \eqref{tra:inverse-color:eq2} in its exponential form and increasing constants gives \eqref{tra:inverse-color:eq1}.
\end{proof}
\subsection{Balanced spectral buckets}
\begin{proposition}\label{tra:inverse-color:recursion}
\label{tra:inverse-color:trace-recursion-statement}
\textbf{Trace bound on splitting the cube.} We next show the following implication when \(d\) is sufficiently large. Split the \(d\) bits contiguously into groups of \(\lfloor d/2\rfloor,\lceil d/2\rceil\), and put \(m,q\) equal to 2 to those respective powers. Suppose for some \(s\ge 2\) that
\[
                              Z_m(s), Z_q(s)\le 2.
\]
Then for \(p=s(1+C_0/\log n)\), with an absolute constant \(C_0\), we have
\begin{equation}
                              Z_n(p)\le \exp(n^{0.60}).                         \label{tra:inverse-color:eq5}
\end{equation}
In particular the loss here has \(\log\) much smaller than \(n\); the order multiplier matters as well, and has been chosen to work with balanced splitting.
\end{proposition}
\begin{proof}\label{tra:inverse-color:recursion-proof}
\label{tra:inverse-color:projection-overlap-and-bucket-summation}
Regard positions as cells in the grid above, with column index coming from the first group of bits. Write \(H\) for the subgroup permuting positions independently within rows, \(J\) that within columns. The full sweep is a product of independent sweeps within rows followed by independent sweeps within columns. Let \(K_H,K_J\) denote the respective group algebra operators, lifted to the full regular representation, so that $T_n=K_JK_H$, with the row sweeps acting first. Products over individual lines in \(H\) or \(J\) correspond within that subgroup's group algebra to tensor products. Use the positive semidefinite operators
\[
                      D_H=K_H K_H^*,\qquad D_J=K_J^* K_J .
\]
Each local factor has the same positive-square moments as the
corresponding child sweep. With these orientations,
Lemma~\ref{lem:positive-power-comparison} gives
\begin{equation}
 Z_n(p)\le\left\|D_H^{p/4}D_J^{p/4}\right\|_4^4.
 \label{tra:inverse-color:eq6}
\end{equation}

For the operator of a row in \(D_H\) divide positive eigenvalues \(\alpha\) into buckets indexed by nonnegative integers \(\ell\), according to
\[
                            \ell\le -s\log\alpha <\ell+1 .
\]
Use its spectral projections onto the buckets, obtained as elements of the row group algebra (e.g. by polynomial calculus in the row regular representation). Do likewise for each column in \(D_J\). For any choice of one nonempty bucket per row and per column, denote the ranks of the projections within their individual regular representations by \(r_z\), where \(z\) ranges over lines (rows and columns). We have
\begin{equation}
                                   r_z\le 2 e^{\ell_z+1}                        \label{tra:inverse-color:eq7}
\end{equation}
by the trace moment hypothesis in the respective smaller cubes. Denote the combined row projection in the big regular representation by \(E\), and the combined column projection by \(F\); projections of individual lines within rows (and likewise within columns) combine by product in the subgroup. All projections and spectral restrictions can be lifted from the subgroup algebras as used here, since restriction of the regular representation to a subgroup acts by copies of its regular representation.

The spectral buckets preserve the child trace budget as a bound on
regular rank. The next calculation converts the probability estimate
into the overlap needed to sum those buckets:
\begin{equation}
               \|E F\|_4^4\le
                  \exp(n^{0.59})\prod_{z} r_z^{\,1+C_1/\log n}                 \label{tra:inverse-color:eq8}
\end{equation}
with an absolute \(C_1\). Apply Lemma~\ref{lem:coefficient-compatibility} to these line
projections. Their squared coefficient masses and the caps of their
normalized squared densities are both $r_z$. Hence
\[
 \|EF\|_4^4\le\frac{n!}{|H||J|}
       \left(\prod_zr_z\right)
       \mathbb P_\rho(\text{reversible in order}),
\]
where $\rho$ is an independent product of line densities bounded by
$r_z$. The lemma's inversion step matches the chronological convention
of Proposition~\ref{tra:inverse-color:compatibility}. Applying that
proposition with $L_R+L_C=\sum_z\log r_z$, and using
\[
 \log\frac{n!}{|H||J|}=n+O(\sqrt n\log n),
\]
proves \eqref{tra:inverse-color:eq8}.

Expand the product in \eqref{tra:inverse-color:eq6} according to the bucket assignments (zero eigenvalues contribute nothing). For one assignment, the term \(D_H^{p/4} E F D_J^{p/4}\) in the product has Schatten 4-norm at most
\[
              \exp\left(-\frac{p}{4s}\sum_z\ell_z\right)\|E F\|_4.
\]
This follows by the ideal property of that norm with the outside operators restricted by the projections (using their operator norms). Use \eqref{tra:inverse-color:eq7},\eqref{tra:inverse-color:eq8} and the triangle inequality. Choose \(C_0>C_1+1\). For an absolute \(C'\), we get
\[
 \begin{split}
           \|D_H^{p/4}D_J^{p/4}\|_4
           &\le \exp(n^{0.59}/4)
              \left[C'\sum_{\ell\ge0}e^{-\ell/(4\log n)}\right]^{q+m}.
 \end{split}
\]
Since \(q+m=O(\sqrt n)\), \eqref{tra:inverse-color:eq6} gives \eqref{tra:inverse-color:eq5} for sufficiently large \(n\), uniformly in \(s\).
\end{proof}
\subsection{An exact endpoint kernel and sparse forests}
\begin{lemma}\label{tra:inverse-color:sparse}
\label{tra:inverse-color:sparse-statement}
The loss in \eqref{tra:inverse-color:eq5} requires a separate bound for sparse first-row levels. Call \(k=n-\lambda_1\) the level of an irreducible, where \(\lambda_1\) is the length of the first row of its diagram. We show, for some absolute \(c_2>0\), all sufficiently large \(n\), and
\[
                           1\le k\le n^{0.64},
\]
that the sweep has squared operator norm (square of its largest singular value), within each level \(k\) irreducible, bounded by
\begin{equation}
                                      n^{-c_2 k}.                              \label{tra:inverse-color:eq9}
\end{equation}
\end{lemma}
\begin{proof}\label{tra:inverse-color:sparse-proof}
\label{tra:inverse-color:endpoint-kernel-and-forest-proof}
Use the permutation action on ordered injective \(k\)-tuples of positions. The level \(k\) irreducibles occur there but are orthogonal to functions missing any of the coordinates, since such functions are in tuple modules for \(k-1\). Indeed by the branching rule (or Young's rule) the tuple module for \(b\) coordinates contains the irreducibles whose first rows have length at least \(n-b\); it is induced from a trivial representation on \(n-b\) positions. We bound the sweep kernel on the part in question in the tuple module.

Consider an entry for input positions \(x\) and output positions \(y\) (ordered tuples). Orientation of the kernel, or taking its transpose, will not matter for the estimates. Paths of these cards through the sweep, if they realize the entry, are uniquely specified by the endpoints, since at stage \(r\) bits \(1,\ldots,r\) have been set to their end values and the other bits have the start values. For two cards \(u,v\) in the tuple let
\begin{itemize}
\item \(b_{uv}\) be the last bit where their starts differ;
\item \(a_{uv}\) be the first bit where their ends differ.

\end{itemize}
The entry probability is
\begin{equation}
          n^{-k} Q_{[k]}(x,y),\qquad
          Q_T(x,y)=\prod_{\{u,v\}\subset T} w(b_{uv},a_{uv}),\quad
 w(b,a)=
 \begin{cases}
 1&a<b,\\
 2&a=b,\\
 0&a>b.
 \end{cases}                                                                  \label{tra:inverse-color:eq10}
\end{equation}
In fact \(a>b\) gives impossible paths (collision), \(a=b\) means they share the switch at that bit and use distinct exits, and in the other case they do not meet at a switch. If no collision occurs, at a switch used by both cards of some pair the probability factor is \(1/2\) rather than the product \(1/4\); other used switches with one card just give the factor \(1/2\). This also explains \eqref{tra:inverse-color:eq10} for any smaller tuple, using the corresponding number of cards instead of \(k\) in the prefactor.

On projecting onto the sum of the level \(k\) irreducibles we can replace \(Q_{[k]}\) by
\[
                    \widetilde Q=\sum_{T\subset [k]} (-1)^{k-|T|} Q_T .
\]
Indeed the added matrices have entries that omit at least one tuple coordinate and vanish between the indicated projected parts. Consider the graph on tuple labels whose edges indicate that the paths specified by \(x,y\) touch a common switch, whether or not they could occur jointly. If this graph has an isolated vertex, \(\widetilde Q=0\), by cancellation. The square of the Hilbert-Schmidt norm of the sweep on the projected part is therefore at most
\begin{equation}
             2^k n^{-2k}
               \sum_{T\subset[k]} \sum_{x,y}
                           Q_T(x,y)^2\, \mathbf 1_{\{\text{no isolates}\}}.   \label{tra:inverse-color:eq11}
\end{equation}

Here are probabilistic bounds for \eqref{tra:inverse-color:eq11}. For fixed \(T\), change measure using one factor \(Q_T\). Up to a total mass factor at most 1, and dropping distinctness constraints on the additional paths for upper bounds, the law now can be sampled as follows:
\begin{itemize}
\item Realize an independent full random sweep on all positions. Pick \(|T|\) distinct start positions uniformly and use their start/end paths for \(T\).
\item For the other labels, use independent paths with starts and ends all uniform independent positions.

\end{itemize}
This follows from \eqref{tra:inverse-color:eq10}: on distinct tuples, \(n^{-2k}Q_T(x,y)\) gives exactly this weight with the factor \((n)_{|T|}/n^{|T|}\), \((n)_a\) denoting the falling factorial. Only the positions within \(T\) are needed to evaluate the remaining factor \(Q_T\), which in this construction equals \(2\) to the power the number of meetings at switches among the paths in \(T\).

Condition on the full sweep realization used in this sampling. We need two bounds.
\paragraph{Meeting count.} The number of meetings among any set of \(b\ge1\) realized full-sweep paths is at most \(b\log_2(b)/2\). To see this split paths by last input bit. All meetings before the last stage respect this split, and the two parts have sweeps on smaller cubes up to then. Meetings at the last stage contribute a matching between the two parts. Induction gives the bound, since for part sizes \(b',b-b'\),
\[
        b\log_2 b-b'\log_2 b'-(b-b')\log_2(b-b')
                            \ \ge\ 2\min(b',b-b')
\]
(with zero summands interpreted by continuity); this follows also by concavity of binary entropy.
\paragraph{Forest probability.} For any specified forest on the \(k\) labels, the probability all its edges indicate common switches is at most
\[
                                    (C d/n)^{\#\text{edges}} .
\]
Indeed sample its paths along the forest in parent-before-child order. If a child is one of the additional independent paths, for each stage its switch (out of \(n/2\)) is uniform, from the uniform endpoints. If it belongs to \(T\), we sample among at least \(n-k\) remaining full-sweep paths, of which at most two at each stage use the switch in question on its parent's path. The estimate follows by testing the edges as the child paths are sampled, since \(k\le n^{0.64}\).

For the event with no isolates, take spanning trees in each component of the common-switch graph. For component sizes \(b_1,\ldots,b_j\ge2\), the first bound above shows that the remaining weight \(Q_T\) is at most \(\prod b_i^{b_i/2}\), since paths from \(T\) in different components cannot meet. For an upper bound on its expectation on the event, sum over spanning forests using this weight for forest component sizes and using the probability bound from the second estimate (we can thereby count outcomes more than once). Forests with \(j\) components number at most \(2^k k^{k-j}\), by choosing roots and parents. Also \(b^{b/2}\le C^b k^{(b-1)/2}\) for \(b\le k\), since \(b^{(b-1)/2}\le k^{(b-1)/2}\) and \(\sqrt b\le C^b\). Thus \eqref{tra:inverse-color:eq11} is bounded by
\begin{equation}
                      C^k
              \sum_{1\le j\le k/2}
                    \left(\frac{C d\, k^{3/2}}{n}\right)^{k-j},               \label{tra:inverse-color:eq12}
\end{equation}
where the constants absorb the subset counts. If there is no such \(j\), \eqref{tra:inverse-color:eq11} is zero. Otherwise \(d=\log_2 n\), \(k^{3/2}\le n^{0.96}\), and \(k-j\ge k/2\); \eqref{tra:inverse-color:eq12} is at most \(n^{-c_2 k}\) with some absolute \(c_2>0\) for sufficiently large \(n\). Since the tuple action applies the group algebra operator within the irreducibles (taking adjoints if using the other kernel orientation), we have proved \eqref{tra:inverse-color:eq9}.
\end{proof}
\begin{lemma}\label{tra:inverse-color:dimension}
\label{tra:inverse-color:dimension-statement}
We specify why this covers the relevant small dimensional part. Any irreducible with diagram height bigger than \(n/2\) is killed by a switch-layer average: the switch subgroup is a Young subgroup with \(n/2\) parts of size 2, and by Young's rule such an irreducible has no invariants there (cannot have a strictly increasing column for a semistandard filling with \(n/2\) symbols). Moreover for all sufficiently large \(n\), any remaining irreducible not of level \(\le n^{0.64}\) has dimension at least
\[
                                 \exp(n^{0.62}).
\]
\end{lemma}
\begin{proof}\label{tra:inverse-color:dimension-proof}
\label{tra:inverse-color:dimension-proof-details}
Here is one quick verification. If both row and column lengths (longest in the diagram) are smaller than \(n/10\), every hook is bounded by \(n/5\), giving the result by the hook-length formula. Otherwise take a direction with length at least \(n/10\), transposing if necessary (dimension unchanged). Outside this long row there are at least \(\lfloor n^{0.64}\rfloor\) boxes by hypothesis (if transposed, by original height at most \(n/2\)). Take a subdiagram with this row and just \(b=\lfloor n^{0.64}\rfloor\) boxes below it, shrinking from corners. Its standard tableaux already give the required dimension lower bound, since they extend to the full diagram. In fact one may fill the first \(b\) places of the row first, then interleave a standard order on the lower \(b\) boxes with the rest of the first row arbitrarily. For large \(n\) this gives at least \(2^{\Omega(b)}\) ways.
\end{proof}
\subsection{The bounded-order induction}
\begin{theorem}\label{tra:inverse-color:moment}
There is an absolute finite $s_*>0$ such that $Z_n(s_*)\le1+1/8$ for every dyadic $n$. 
\end{theorem}
\begin{proof}[Proof of Theorem~\ref{tra:inverse-color:moment}]
\label{tra:inverse-color:bounded-order-induction-and-completion}
Choose a fixed $P\ge2$ with $Pc_2\ge4$, where $c_2$ is the
constant in \eqref{tra:inverse-color:eq9}. At first-row level
$1\le k\le n^{.64}$ there are at most $2^k$ shapes, each of dimension
at most $n^k$. Thus for every $u\ge P$, their contribution to $Z_n(u)$
is at most
\[
 \sum_{k=1}^{\lfloor n^{.64}\rfloor}
       2^k n^{2k-u c_2k}=o(1).
\]
Lemma~\ref{tra:inverse-color:dimension} annihilates shapes of height
more than $n/2$ and gives dimension at least $\exp(n^{.62})$ for all
remaining shapes outside this sparse range.

Apply Lemma~\ref{lem:balanced-moment-closure} with positive-square
convention $\nu=2$ and $\varepsilon=1/8$. The sparse and annihilated
shapes form its first class and contribute at most $1/16$ for large
$n$, uniformly above $P$. Proposition~\ref{tra:inverse-color:recursion}
supplies the rough bound from child moments at most $1+1/8<2$, with
\[
 E_n=n^{.60},\qquad H_n=n^{.62},\qquad
 \alpha_d=1+C_0/\log n,\qquad\delta_d=1/\log n.
\]
The required attenuation follows from
\[
 (1+\delta_d)E_n-\delta_dH_n
 =(1+1/\log n)n^{.60}-n^{.62}/\log n\longrightarrow-\infty.
\]
All size thresholds are independent of the inherited exponent, and
$\alpha_d(1+\delta_d)=1+O(1/d)$. The common lemma therefore gives a
finite $s_*$ such that $Z_n(s_*)\le1+1/8$ for all dyadic $n$.

This is a positive-square moment, with matching singular exponent
$2s_*$. Proposition~\ref{e:nonnormal-moment-conversion} gives
chi-square divergence at most $1/8$ after every integer number
$l\ge s_*$ of forward sweeps.
\end{proof}
